% Section 6 -- the withdrawn finding, compressed to a negative-result
% paragraph. Framing kept: the error was OURS. The concession that the
% attribute may behave exactly as documented is now a pointer to Appendix I
% (app:daglayer), not stated here. Never quote the withdrawn depth figures.

\section{A withdrawn finding: \texttt{dag\_layer} and the catching check}
\label{sec:dag-layer}

The graph release publishes a per-declaration attribute, \texttt{dag\_layer},
which we took to be distance from primitives. We built a depth metric on it,
measured a machine--human gap, and promoted that gap to a stated
counter-finding. The column is not a depth measure as we used it: layer zero
holds \layerZeroNodes{} declarations whose maximum in-degree is zero (uncited
leaves, not axioms), the highest layer holds \texttt{Eq.refl}, \texttt{Set}
and \texttt{LE.le}, and the column correlates \layerCorrLogIndegree{} with log
in-degree, so a high value means heavily cited infrastructure and the scale runs
opposite to the direction we read into it. The error was ours, and
Appendix~\ref{app:daglayer} gives the reading of the column under which it
behaves as documented, together with a second hazard for other consumers of the
release. What we measured is therefore a centrality proxy, and a
difference in it is a statement about citation weight rather than derivation
length. That is the inference \citet{li2026network} rule out in their
Finding~3, which we cite in Section~\ref{sec:background} and then walked into
on a different column of the same release. We withdraw the metric rather than
re-pointing it at a claim about centrality that it would support, because
choosing a quantity's meaning after seeing its comparison is a researcher
degree of freedom. The check that would have caught it is cheap enough to run
on any released structural attribute before building on it: correlate the
attribute against in-degree, then inspect its extremes and ask whether the
declarations at each end match the meaning its name implies.
